\section{Capítulo~\ref{sec:serocomputer}. Neuronas \nt{SERO}}

Las propiedades de las propias neuronas \nt{SERO} (ritmo, autoinhibición, signo según el receptor, subsistemas) están medidas directamente; los cálculos del capítulo (regulador, filtro, lógica) se siguen de sus ecuaciones; $\tau_c$, el coeficiente de difusión de la serotonina y el número de sitios de liberación son estimaciones; el papel del lazo como regulador de nivel en el cerebro es una hipótesis (en el núcleo del rafe la autoinhibición es puntual y solo da pausas breves), y el papel de \nt{SERO} como horizonte es una hipótesis con experimentos de conducta a su favor.

{\footnotesize
\setlength{\tabcolsep}{3pt}\renewcommand{\arraystretch}{1.15}
\begin{longtable}{>{\raggedright}p{0.5cm}>{\raggedright}p{4.0cm}>{\raggedright}p{5.6cm}>{\raggedright}p{2.3cm}>{\raggedright\arraybackslash}p{1.7cm}}
\toprule
N.º & Proposición & Experimento y qué se midió & Fuente & Estado \\
\midrule\endfirsthead
\toprule
N.º & Proposición & Experimento y qué se midió & Fuente & Estado \\
\midrule\endhead
1 & El signo de la acción de la serotonina lo elige el receptor del destinatario (proposición~\ref{prop:receptor}) & Los receptores de serotonina se dividen en familias según su farmacología y su mecanismo: 5-HT$_{1A}$ abre canales de potasio a través de una proteína G e inhibe la célula; 5-HT$_{2A}$ y el ionotrópico 5-HT$_3$ excitan. Un mismo neurotransmisor da respuestas opuestas en células distintas. & \cite{BarnesSharp1999} & medido \\
2 & La neurona \nt{SERO} en vigilia emite de forma regular varias espigas por segundo & Registro de neuronas del núcleo dorsal del rafe en gatos en libertad de movimiento: descarga rítmica lenta de $2.8\pm0.2$ espigas/s en vigilia tranquila; en el sueño lento la frecuencia cae un $34$--$68\%$, en el sueño REM, un $98\%$. En ratas anestesiadas, $1.7\pm0.5$ espigas/s, muy regular. & \cite{TrulsonJacobs1979, GagnonParent2014, JacobsAzmitia1992} & medido \\
3 & La neurona \nt{SERO} es un marcapasos: no necesita un empujón para cada espiga, pero sí un aporte de fondo constante (en la rebanada, noradrenalina a través de receptores $\alpha_1$; en el organismo, desde \nt{NORA}) & En la rebanada de cerebro las células descargan de forma lenta y regular; tras cada espiga hay una poshiperpolarización de $200$--$800\ms$. En la rebanada hay menos células activas espontáneamente que en el organismo: el ritmo regular necesita una entrada tónica de noradrenalina a través de receptores $\alpha_1$, cuyo bloqueo lo apaga. & \cite{VandermaelenAghajanian1983} & medido \\
4 & Casi todos los receptores son metabotrópicos; la respuesta es lenta: sube y baja en cerca de un segundo & Todas las familias de receptores salvo 5-HT$_3$ están acopladas a proteínas G. En la rebanada, la liberación evocada de serotonina produce a través de 5-HT$_{1A}$ una corriente inhibidora que sube y baja en un segundo. & \cite{BarnesSharp1999, CourtneyFord2016} & medido \\
5 & En las regiones de destino la serotonina sale al volumen, no solo a la hendidura; en el propio núcleo del rafe la inhibición de las vecinas es punto a punto & Voltamperometría cíclica rápida en rebanadas: la liberación por espiga es igual en una región con sinapsis y en el núcleo del rafe, donde casi no hay sinapsis; no depende del bloqueo de la recaptación ni de los receptores; hay muchas menos moléculas por terminal que transportadores y receptores, y la concentración extracelular coincide con la afinidad del receptor principal. En el núcleo del rafe la inhibición a través de 5-HT$_{1A}$ es punto a punto: el transportador impide que la serotonina se acumule fuera de la hendidura. & \cite{Bunin1998, CourtneyFord2016} & medido \\
6 & La concentración según~\eqref{eq:seroneuron} decae por la recaptación; $\tau_c$ del orden de un segundo es una estimación & Voltamperometría en el cerebro vivo: la serotonina extracelular la limitan ante todo la recaptación y la degradación, no la síntesis. En los trabajos citados no hay medida directa de $\tau_c$; el orden de magnitud es una estimación del capítulo. & \cite{Hashemi2012serotonin} & medido; $\tau_c$, estimación \\
7 & NO y NO-O con un solo marcapasos; completitud de la lógica \nt{SERO} (proposición~\ref{prop:serocomplete}, fig.~\ref{fig:serogates}) & Se sigue de~\eqref{eq:seroneuron}: un marcapasos inhibible es un NO-O, y el NO-O es funcionalmente completo. No hay experimentos en que se haya construido lógica con neuronas \nt{SERO}; la condición de volúmenes separados no se cumple en el cerebro. & deducción en el capítulo & teoría \\
8 & La serotonina inhibe a la propia neurona \nt{SERO} a través de receptores situados en ella & En ratas anestesiadas, los agonistas de 5-HT$_{1A}$ por vía intravenosa y por iontoforesis inhiben la descarga de las neuronas del rafe dorsal con la misma relación dosis-respuesta que la propia serotonina; los agonistas de 5-HT$_{1B}$ casi no actúan. En la rebanada, la serotonina endógena de las células vecinas produce a través de 5-HT$_{1A}$ una corriente y una pausa breve en la descarga, sin cambiar la frecuencia media. & \cite{Sprouse1987, CourtneyFord2016} & medido \\
9 & En el modelo, el lazo a través de un volumen común es un regulador de nivel: la dispersión y la constante de tiempo se reducen $1+G$ veces~\eqref{eq:feedback}; que el lazo funcione así en el cerebro es una hipótesis & Fórmula clásica del amplificador con realimentación negativa; en el capítulo se deduce de nuevo. La ganancia de lazo $G$ del sistema \nt{SERO} no se ha medido. Medido: en la rebanada la autoinhibición es punto a punto, da una pausa breve y no cambia la frecuencia media. & \cite{Black1934feedback, CourtneyFord2016} & teoría; en el cerebro, hipótesis \\
10 & La concentración es una media móvil de la frecuencia, un filtro paso bajo~\eqref{eq:lowpass}, fig.~\ref{fig:lowpass} & Solución de la ecuación lineal~\eqref{eq:seroneuron}; la figura es un cálculo con esta fórmula para $\tau_c=1\,\text{s}$. No hay un modelo aparte en \texttt{models/}. & deducción en el capítulo & teoría \\
11 & La tortuosidad de los espacios frena la difusión de una molécula pequeña unas $2.6$ veces & Método de fuente puntual: iontoforesis de tetrametilamonio y registro de su concentración con un electrodo ionoselectivo en rebanadas y en el cerebro vivo. Tortuosidad $\lambda\approx1.6$, es decir, el $D$ efectivo es $\lambda^2\approx2.6$ veces menor que el libre; fracción de volumen extracelular, cerca de $0.2$. El coeficiente de difusión de la propia serotonina en el tejido no se midió en estos trabajos; en el capítulo es una estimación. & \cite{Sykova2008} & medido \\
12 & Longitud de un “cable” independiente $\ell\sim\sqrt{D\tau}\approx20\um$~\eqref{eq:diffusionlength}; el número de cables está limitado por el número de tales regiones (proposición~\ref{prop:volumewires}) & Ley de difusión y sustitución de las estimaciones de $D$ y $\tau$ (filas~6 y~11); la proposición~\ref{prop:volumewires} es una consecuencia. No hay experimentos directos que cuenten los canales independientes de la transmisión de volumen. & deducción en el capítulo & teoría \\
13 & La salida de una neurona \nt{SERO} se reparte por enormes zonas del cerebro; se estiman $\sim10^5$ sitios de liberación & Reconstrucción completa de neuronas individuales del rafe dorsal en la rata: todos los axones ascendentes se ramifican mucho, la longitud axonal total de una célula llega a $18.7$~cm, y las ramas de una misma célula alcanzan el estriado, la corteza prefrontal y la amígdala. El número de sitios de liberación por célula no se ha contado directamente. & \cite{GagnonParent2014} & medido; $10^5$, estimación \\
14 & Las neuronas \nt{SERO} se dividen en varios subsistemas con salidas y respuestas distintas & Marcaje viral-genético en ratones: las neuronas que van a la amígdala y a la corteza frontal casi no se solapan en su ramificación, reciben entradas distintas y responden de forma opuesta a un estímulo desagradable; activar y desactivar cada grupo cambia la conducta de manera diferente. & \cite{Ren2018} & medido \\
15 & Condición de paciencia $D<\tau_\gamma\ln(R/r_0)$~\eqref{eq:patience} con descuento exponencial; con el medido, hiperbólico, $D<\tau_\gamma(R/r_0-1)$ & Se sigue del descuento $e^{-D/\tau_\gamma}$ o $1/(1+D/\tau_\gamma)$; deducción en el capítulo. Monos, elección con esperas de segundos: el descuento se parece más a una hipérbola que a una exponencial; la respuesta de las neuronas \nt{DOPA} a la señal que anuncia la recompensa diferida cae con la espera también como una hipérbola. & deducción en el capítulo; \cite{KobayashiSchultz2008} & teoría \\
16 & \nt{SERO} fija el horizonte $\tau_\gamma$ (sección~\ref{sec:horizon}) & La activación optogenética de las neuronas \nt{SERO} del rafe dorsal en ratones alarga la espera de una recompensa diferida y aumenta la persistencia en buscar una recompensa que se ha vuelto escasa. En humanos, bajar la serotonina (dieta sin triptófano) hace más empinado el descuento de la recompensa diferida. Cómo se convierte la concentración en $\tau_\gamma$ no se sabe. & \cite{Doya2002, Miyazaki2014, Lottem2018, Schweighofer2008} & hipótesis \\
\bottomrule
\end{longtable}}
